% =========================================================================
% Compilation draft — clean single-column layout for review and PDF preview.
% Final submission will use the Springer Nature sn-jnl template.
%
% Practical Limits and Numerical Reliability in
% Exact Diagonalization of Quantum Spin Chains
%
% All timings from a single logged calibration session on 2026-10-07:
%   session 20261007_070101
% Every numerical value in this manuscript is reproducible from the
% archived session CSVs with the accompanying reproduce.py script.
% =========================================================================

\documentclass[pdflatex,sn-mathphys-num]{sn-jnl}

\usepackage{amsmath,amssymb}
\usepackage{graphicx}
\usepackage{booktabs}
\usepackage{url}

\graphicspath{{figures/}}
\newcommand{\safeincludegraphics}[2][]{%
  \IfFileExists{#2}{\includegraphics[#1]{#2}}{%
    \IfFileExists{figures/#2}{\includegraphics[#1]{figures/#2}}{%
      \IfFileExists{results/figures/#2}{\includegraphics[#1]{results/figures/#2}}{%
        \fbox{\parbox[c][0.24\columnwidth][c]{0.78\columnwidth}{%
          \centering Missing generated figure: \texttt{#2}}}}}}}

\setlength{\parskip}{0.5ex}
\setlength{\emergencystretch}{2em}

\begin{document}

\title{Practical Limits and Numerical Reliability in Exact Diagonalization of Quantum Spin Chains}

\author*[1]{\fnm{Saleh} \sur{Swaid}}\email{sswaid991@gmail.com}
\author[2]{\fnm{M.~Imran} \sur{Jamil}}

\affil*[1]{\orgdiv{Department of Physics}, \orgname{Amran University},
\orgaddress{\city{Amran}, \country{Yemen}}}
\affil[2]{\orgdiv{Department of Physics}, \orgname{University of
Management and Technology}, \orgaddress{\city{Lahore},
\country{Pakistan}}}

\abstract{%
\noindent
Exact diagonalization (ED) remains a reference tool for small quantum
spin chains, yet several practical choices are made by convention
rather than by measurement. Using a single-threaded NumPy/SciPy
implementation on one laptop, we benchmark three such choices for
open spin-$1/2$ chains. First, we compare dense and sparse Lanczos
workflows for obtaining the ground-state energy. The sparse solver
becomes faster at a Hilbert-space dimension $D \approx 2\times10^{2}$,
with workflow speedups of $\approx 3.4\times10^{2}$ at $D \approx 2000$
and $\approx 1.1\times10^{3}$ at $D = 4096$ when the dense reference is
\texttt{numpy.linalg.eigvalsh}. Replacing the dense reference with
\texttt{scipy.linalg.eigvalsh} reduces the ratio to
$\approx 6.6\times10^{2}$ at $D = 4096$, because the two libraries
differ by a factor of $2.3$ on the same matrix. Restricting either
dense call to the lowest eigenvalue via \texttt{subset\_by\_index=[0,0]}
provides an additional $\approx 2$--$13\%$ saving at
$D \geq 924$; at $D \leq 512$ the saving reaches $\approx 30\%$ but
lies within timing noise. Second, the leading finite-size correction
to $E_0/N$ in open chains is a $1/N$ surface term: for the critical
transverse-field Ising chain, a three-parameter fit recovers the exact
bulk energy to $\approx 10^{-4}$ at $N_{\max} = 16$, and a
four-parameter fit including the $N^{-3}$ term reduces the
extrapolation error to $\approx 6.3\times10^{-6}$. Restricting to the
$P = +1$ parity sector gives $1.9$--$3.8\times$ speedups. Third,
half-chain entanglement entropies of the Heisenberg chain show a
pronounced $N \bmod 4$ dependence; single-branch fits on three or
four points give effective central charges that differ by a factor of
three between the two branches. Data and scripts are openly
available.
}

\keywords{Exact diagonalization; Quantum spin chains; Lanczos method;
Finite-size scaling; Symmetry reduction; Entanglement entropy}

\maketitle

% =========================================================================
\section{Introduction}
\label{sec:intro}

Exact diagonalization (ED) is the reference method for small quantum
many-body systems. It produces eigenvalues and eigenvectors of a
finite-size Hamiltonian without variational or stochastic
approximation, and serves as a validation baseline for tensor-network
and quantum-Monte-Carlo codes. Its limitation is the exponential
growth of the Hilbert-space dimension. General introductions are
given in Sandvik's review~\cite{Sandvik2010}; mature packages such as
QuSpin~\cite{Weinberg2017} exist, and tensor-network alternatives are
reviewed in Ref.~\cite{Schollwock2011}.

Despite this maturity, several practical questions inside an ED
implementation are often answered by convention rather than by
measurement. We focus on three.

\emph{Dense versus sparse eigensolvers.} Dense diagonalization via
\texttt{numpy.linalg.eigvalsh} (LAPACK) computes the full spectrum at
cost $\mathcal{O}(D^3)$ in time and $\mathcal{O}(D^2)$ in memory. The
implicitly restarted Lanczos method of ARPACK, exposed through
\texttt{scipy.sparse.linalg.eigsh}~\cite{Lehoucq1998}, computes only
$k$ extremal eigenvalues at cost $K$ matrix-vector products of
$\mathcal{O}(\mathrm{nnz})$ each. The two are not, strictly speaking,
interchangeable---one returns the entire spectrum, whereas the other
targets only a small number of extremal eigenvalues. Nevertheless,
in practical ED workflows a dense diagonalization is sometimes used
when only the ground-state energy is ultimately required. We
therefore benchmark this workflow explicitly, while treating the
resulting dense-to-sparse ratios as workflow speedups rather than as
intrinsic solver speed ratios.

\emph{Finite-size extrapolation.} The appropriate fit variable
depends on boundary conditions. For open chains the chain ends
contribute a surface energy that makes the leading correction to
$E_0/N$ of order $1/N$; for periodic chains the leading correction is
of order $1/N^{2}$. The two ans\"atze are not interchangeable.

\emph{Symmetry reduction.} Its benefit is obvious in principle but
rarely quantified as a function of system size, nor separated from
its effect on solver convergence.

This is a benchmarking study. Its contribution is to replace
inherited conventions by measured values, with their uncertainties
and limitations stated explicitly. The implementation, raw data, and
figure scripts are available at
\url{https://github.com/sswaid991-bit/ed-spin-chains}.

% =========================================================================
\section{Models and methods}
\label{sec:methods}

\subsection{Models and conventions}
\label{sec:models}

We consider two spin-$1/2$ chains with open boundary conditions
(OBC). All Hamiltonians use Pauli matrices.

\paragraph{Transverse-field Ising (TFI) chain.}
\begin{equation}
H_{\mathrm{TFI}} = -J \sum_{i=1}^{N-1} \sigma^z_i \sigma^z_{i+1}
                  - h \sum_{i=1}^{N} \sigma^x_i ,
\label{eq:tfi}
\end{equation}
at the critical point $h = J = 1$. The model conserves neither $S^z$
nor $S^x$, so the full $2^N$-dimensional space is used unless the
parity sector is invoked. The bulk energy per site is
$E_0/N \to -4/\pi$; the open-chain ground-state energy has the closed
form~\cite{Pfeuty1970}
\begin{equation}
E_0(N) = 1 - \frac{1}{\sin\!\left(\frac{\pi}{4N+2}\right)},
\label{eq:tfi-closed}
\end{equation}
which we use as an independent check throughout.

\paragraph{Heisenberg antiferromagnet.}
\begin{equation}
H_{\mathrm{XXZ}} = \sum_{i=1}^{N-1}
  \left( \sigma^x_i \sigma^x_{i+1}
       + \sigma^y_i \sigma^y_{i+1}
       + \Delta\, \sigma^z_i \sigma^z_{i+1} \right),
\label{eq:xxz}
\end{equation}
at $\Delta = 1$, in the $S^z=0$ sector of dimension
$D_{S^z} = \binom{N}{N/2}$. In the Pauli convention
\begin{equation}
H = \sum_{i=1}^{N-1} \vec\sigma_i \cdot \vec\sigma_{i+1}
  = 4\sum_{i=1}^{N-1} \vec S_i \cdot \vec S_{i+1},
\label{eq:heis-convention}
\end{equation}
so the bulk energy per site is
\begin{equation}
\epsilon_\infty^{\mathrm{Heis}} = 1 - 4\ln 2 \approx -1.77259,
\label{eq:heis-limit}
\end{equation}
four times the more familiar $-(\ln 2 - 1/4) \approx -0.4431$ of the
spin-exchange convention~\cite{Hulthen1938}. Data for $\Delta = 0.5$
and $2.0$ are given in Appendix~\ref{app:xxz}.

\subsection{Implementation}
\label{sec:impl}

The implementation is a single Python module of under 600 lines using
only NumPy and SciPy. Basis states are integers
$s = \sum_i s_i 2^i$ with $s_i \in \{0,1\}$. The $S^z=0$ sector
consists of all states with $N/2$ up spins. The Hamiltonian is built
in Compressed Sparse Row (CSR) format. The stored number of nonzeros
was verified against closed forms:
$\mathrm{nnz} = (N+1)2^N$ for the TFI chain and
$\mathrm{nnz} = (N/2+1) D_{S^z}$ for the Heisenberg chain, exact in
all six cases tested ($N=4,6,8,10,12,14$ for Heisenberg).

Dense diagonalization uses \texttt{numpy.linalg.eigvalsh}
(eigenvalues only). The sparse solver is the implicitly restarted
Lanczos method implemented by \texttt{scipy.sparse.linalg.eigsh}
through ARPACK, with $k=1$, \texttt{which='SA'}, a fixed-seed normal
starting vector \texttt{v0}, and \texttt{tol=0} (which instructs
ARPACK to use the machine-precision bound on the residual) unless
stated otherwise. A fixed \texttt{v0} makes the number of Lanczos
iterations reproducible. Matrix-vector products are counted by
wrapping the matrix in a \texttt{LinearOperator}.

\subsection{Timing protocol}
\label{sec:timing-protocol}

All timings are wall-clock \texttt{time.perf\_counter} measurements
of the solver call only (matrix construction excluded) and are
reported as the \emph{minimum} over $n$ repetitions, following the
\texttt{timeit} convention, with $n=20$ for $D \le 1024$, $n=10$ for
$1024 < D \le 2048$, and $n=5$ for $D > 2048$. One untimed warm-up
call precedes the timed repetitions. The eigenvalues, sector
dimensions, and Lanczos iteration counts reported here are
reproducible to within floating-point rounding, given the same
library versions and starting vector; wall-clock times depend on the
machine and on the session.

The environment is a single CPU core on a Windows~11 laptop
(Intel Core i7-class processor, 4 physical cores, 16~GB RAM, AVX2
and FMA3 but no AVX512) with Python~3.11, NumPy~1.24.3, and
SciPy~1.15.3. Both OpenBLAS thread pools loaded at runtime (NumPy's
\texttt{libopenblas64} and SciPy's \texttt{libscipy\_openblas}) were
confirmed with \texttt{threadpoolctl} to run with a single thread;
because environment-variable conventions are not honored uniformly
across OpenBLAS builds, we additionally enforce single-threaded
execution at runtime via \texttt{threadpoolctl.threadpool\_limits(1)}.
The exact library files and OpenBLAS versions are recorded in the
session log archived with this paper. We make no claim about other
hardware: the numerical ratios are specific to this configuration,
and the qualitative conclusions are the only part we expect to be
general.

All timing tables in this paper cite a single logged calibration
session on the same machine, with single-threaded BLAS enforced.
Within that session, repeated invocations of the same dense
computation at $D=4096$ are reproducible to $1$--$3\%$ and repeated
invocations of the sparse solver to $10$--$25\%$. The larger
relative scatter in the sparse timings reflects their small absolute
magnitude ($\sim 20$~ms at $D=4096$), where Python and ARPACK
overhead is a significant fraction of the total. These figures
characterize short-term reproducibility within one measurement
session; variation across sessions, thermal loads, and background
processes was not characterized and may be larger.

\subsection{Dense solvers and library dependence}
\label{sec:libraries}

Two numerically equivalent dense eigenvalue routines are in common
use: \texttt{numpy.linalg.eigvalsh} and \texttt{scipy.linalg.eigvalsh}.
Both bind to LAPACK but link against different OpenBLAS builds
(see the \texttt{threadpool\_info} entries in the session log
archived with this paper for the exact library files). On the same
matrices, with single-threaded BLAS, we measured:

\begin{itemize}
  \item At $D \geq 1024$, \texttt{scipy.linalg.eigvalsh} is
        $2.0$--$2.5$ times faster than \texttt{numpy.linalg.eigvalsh}
        (Table~\ref{tab:fair}), with the ratio reaching $2.26$ at
        $D = 4096$ for the TFI chain and $2.43$ at $D = 3432$ for the
        Heisenberg chain.
  \item Below $D \approx 512$ the two calls are indistinguishable.
  \item Restricting either call to the lowest eigenvalue via
        \texttt{subset\_by\_index=[0,0]} (available only through
        \texttt{scipy.linalg.eigvalsh}, driver \texttt{evr}) provides
        an additional saving of $\approx 2$--$13\%$ at $D \geq 924$.
        At $D \leq 512$ the saving reaches $\approx 30\%$, but the
        absolute times there are a few milliseconds and the
        percentage lies within the timing noise. The saving is small
        at large $D$ because the tridiagonalization step, which
        dominates the $\mathcal{O}(D^3)$ cost, is not avoided by the
        subset request.
\end{itemize}

The dense workflow in this paper is \texttt{numpy.linalg.eigvalsh},
which is the default choice in most ED codes. Users of
\texttt{scipy.linalg.eigvalsh} should expect dense-to-sparse ratios
roughly a factor of $2.3$ smaller at large $D$ than those reported
here; the qualitative preference for sparse Lanczos is unaffected.
We have not isolated the cause of the numpy/scipy difference; it is
consistent with the two packages linking against OpenBLAS builds
compiled with different optimization profiles, but this was not
verified experimentally.

\subsection{Fairness of the dense--sparse comparison}
\label{sec:fairness}

A caveat applies to every ratio in this paper. Dense
\texttt{eigvalsh} computes \emph{all} eigenvalues while sparse
\texttt{eigsh} computes only the ground state, so the two solvers are
answering different questions. This asymmetry does not bias the
comparison appreciably, as shown in Table~\ref{tab:fair}:
restricting the dense computation to the lowest eigenvalue via
LAPACK's MRRR driver reduces the dense runtime by $\approx 2$--$13\%$
at $D \geq 924$ (and up to $\approx 30\%$ at $D \leq 512$, within
timing noise), and does not change the qualitative preference for
sparse Lanczos at any size tested. This is consistent with the
tridiagonalization step, which dominates the $\mathcal{O}(D^{3})$
cost, being unavoidable regardless of how many eigenvalues are
requested. The reported speedups are therefore not an artifact of
the dense solver computing more eigenvalues than necessary.

% =========================================================================
\section{Results}
\label{sec:results}

\subsection{Dense versus sparse diagonalization}
\label{sec:timing}

Table~\ref{tab:timing} lists minimum wall-clock times for the numpy
dense baseline. Below $D \approx 70$ both solvers finish in about a
millisecond and sparse is slower (ratio $0.1$--$0.3$), because fixed
ARPACK and Python overhead dominates. At $D = 252$--$256$ sparse is
already $2.1$--$2.4$ times faster. Interpolating $\ln(\text{ratio})$
against $\ln D$ between the bracketing points gives a break-even at
$D \approx 1.9\times10^{2}$ (TFI) and $D \approx 1.5\times10^{2}$
(Heisenberg), which we quote as $D \approx 2\times10^{2}$---a
working midpoint rather than a precise determination
(Fig.~\ref{fig:crossover}). A factor of ten in favor of sparse is
reached near $D \approx 5\times10^{2}$, and the ratio is
$\approx 4\times10^{1}$ at $D \approx 10^{3}$.

\begin{table}[t]
\centering
\small
\caption{Minimum wall-clock time (s) for the ground-state energy.
Dense: \texttt{numpy.linalg.eigvalsh}; sparse:
\texttt{scipy.sparse.linalg.eigsh}, $k=1$, \texttt{which='SA'},
\texttt{tol=0}, fixed-seed \texttt{v0}. Repetitions: $n=20$ for
$D \le 1024$, $n=10$ for $1024 < D \le 2048$, and $n=5$ for
$D > 2048$. A dash means the dense calculation was not attempted.
The ratio is $t_{\mathrm{dense}}/t_{\mathrm{sparse}}$. Values below
$0.001$~s are dominated by the system clock resolution. All values
from session \texttt{20261007\_070101}.}
\label{tab:timing}
\begin{tabular}{lrrrrr}
\toprule
Model & $N$ & $D$ & $t_{\text{dense}}$ & $t_{\text{sparse}}$ & ratio \\
\midrule
\multicolumn{6}{l}{\textit{Heisenberg ($S^z = 0$ sector)}} \\
Heis & 4  & 6     & $<0.001$ & $<0.001$ & --- \\
Heis & 6  & 20    & $<0.001$ & $<0.001$ & --- \\
Heis & 8  & 70    & $<0.001$ & $0.0011$ & $0.3$ \\
Heis & 10 & 252   & $0.0039$ & $0.0017$ & $2.4$ \\
Heis & 12 & 924   & $0.142$  & $0.0033$ & $4.3\times10^{1}$ \\
Heis & 14 & 3432  & $15.29$  & $0.0152$ & $1.0\times10^{3}$ \\
Heis & 16 & 12870 & ---      & $0.0581$ & --- \\
\multicolumn{6}{l}{\textit{Transverse-field Ising (full $2^N$ space)}} \\
TFI  & 4  & 16    & $<0.001$ & $<0.001$ & --- \\
TFI  & 6  & 64    & $<0.001$ & $0.0019$ & $0.2$ \\
TFI  & 8  & 256   & $0.0072$ & $0.0033$ & $2.2$ \\
TFI  & 10 & 1024  & $0.289$  & $0.0055$ & $5.2\times10^{1}$ \\
TFI  & 11 & 2048  & $3.26$   & $0.0095$ & $3.4\times10^{2}$ \\
TFI  & 12 & 4096  & $27.55$  & $0.0240$ & $1.1\times10^{3}$ \\
TFI  & 14 & 16384 & ---      & $0.132$  & --- \\
TFI  & 16 & 65536 & ---      & $0.818$  & --- \\
\bottomrule
\end{tabular}
\end{table}

\begin{figure}[t]
\centering
\safeincludegraphics[width=0.85\textwidth]{Fig1_crossover.pdf}
\caption{Minimum wall-clock time for the ground-state calculation as a
function of Hilbert-space dimension $D$. The break-even occurs near
$D \approx 2\times10^{2}$. The dense and sparse solvers are not
strictly equivalent: dense returns the full spectrum while sparse
returns only the ground state (see Sec.~\ref{sec:fairness}).}
\label{fig:crossover}
\end{figure}

\paragraph{Dense solvers across libraries.}
Table~\ref{tab:fair} lists the same dense calculation performed by
three workflows: \texttt{numpy.linalg.eigvalsh} (full spectrum),
\texttt{scipy.linalg.eigvalsh} (full spectrum), and
\texttt{scipy.linalg.eigvalsh} with \texttt{subset\_by\_index=[0,0]}.
The difference between the two libraries is a factor of
$2.0$--$2.5$ at $D \geq 1024$; the additional saving from restricting
to the lowest eigenvalue is $\approx 2$--$13\%$ at $D \geq 924$, and
up to $\approx 30\%$ at smaller $D$ where the absolute times are in
the millisecond range and the percentages lie within timing noise.

\begin{table}[t]
\centering
\small
\caption{Dense solvers on the same matrices, single-threaded BLAS,
minimum of $n$ repetitions after one untimed warm-up. All values from
session \texttt{20261007\_070101}. All times in seconds.}
\label{tab:fair}
\begin{tabular}{lrrrrr}
\toprule
$N$ & $D$ & $t_{\text{numpy}}$ & $t_{\text{scipy}}$
    & $t_{\text{scipy,subset}}$ & $t_{\text{numpy}}/t_{\text{scipy}}$ \\
\midrule
\multicolumn{6}{l}{\textit{Transverse-field Ising}} \\
8  & 256   & $0.0057$ & $0.0045$ & $0.0023$ & $1.27$ \\
10 & 1024  & $0.220$  & $0.107$  & $0.101$  & $2.06$ \\
12 & 4096  & $27.86$  & $12.32$  & $12.10$  & $2.26$ \\
\multicolumn{6}{l}{\textit{Heisenberg ($S^z = 0$)}} \\
10 & 252   & $0.0046$ & $0.0039$ & $0.0028$ & $1.18$ \\
12 & 924   & $0.118$  & $0.0718$ & $0.0589$ & $1.64$ \\
14 & 3432  & $15.27$  & $6.274$  & $6.126$  & $2.43$ \\
\bottomrule
\end{tabular}
\end{table}

\paragraph{Cost model and measured iteration counts.}
Dense \texttt{eigvalsh} is dominated by tridiagonalization,
$\approx\tfrac43 D^3$ floating-point operations: $9.2\times10^{10}$
at $D = 4096$, which in $27.6$~s corresponds to $\approx 3.3$~GFLOP/s
on one core. For sparse Lanczos the cost is $K$ matrix-vector
products of $\mathcal{O}(\mathrm{nnz})$ each, plus the
reorthogonalization work of ARPACK, so the cost scales as
$K \cdot \mathrm{nnz} \propto K N D$. Table~\ref{tab:matvec} gives
the measured $K$ for the TFI chain.

A commonly assumed value $K \approx 30$ underestimates the measured
counts by $2$--$5\times$; the measured values increase with $N$ over
the range tested. The observed growth is compatible with the
heuristic estimate $K \sim \sqrt{W/\Delta}$ given a bandwidth
$W \propto N$ and a gap $\Delta \propto 1/N$, but we do not claim
this relation as a fitted scaling law; the measured $K$ values
increase in steps of about $10$ because ARPACK checks convergence
only at restarts. Relaxing \texttt{tol} from its tightest setting to
$10^{-10}$ removes $\approx 20$--$40$ iterations and shortens the
runtime by $20$--$40\%$ at every size.

\begin{table}[t]
\centering
\small
\caption{Lanczos matrix-vector counts $K$ and minimum solver times
(s) for the critical TFI chain at two tolerances. The last column is
$t_{\min}/(K\,\mathrm{nnz})$ in nanoseconds. $K$ values are
deterministic given the fixed \texttt{v0} and reproduce exactly
across sessions. All values from session \texttt{20261007\_070101}.}
\label{tab:matvec}
\begin{tabular}{rrrrrrrr}
\toprule
$N$ & $D$ & nnz & $K(0)$ & $K(10^{-10})$ & $t(0)$ & $t(10^{-10})$ & ns/($K$\,nnz) \\
\midrule
8  & 256   & 2304    & 71  & 51  & $0.0026$ & $0.0018$ & $16.1$ \\
10 & 1024  & 11264   & 91  & 71  & $0.0080$ & $0.0047$ & $7.8$ \\
12 & 4096  & 53248   & 101 & 71  & $0.0202$ & $0.0134$ & $3.8$ \\
14 & 16384 & 245760  & 121 & 91  & $0.0865$ & $0.0697$ & $2.9$ \\
16 & 65536 & 1114112 & 141 & 101 & $0.7836$ & $0.6054$ & $5.0$ \\
\bottomrule
\end{tabular}
\end{table}

With $K = 101$ and $\mathrm{nnz} = 5.3\times10^{4}$ at $D = 4096$,
the nominal dense tridiagonalization count is approximately
$9.2\times10^{10}$ floating-point operations, whereas the sparse
matvec count is only of order $K\,\mathrm{nnz} \simeq 5.4\times10^{6}$
sparse matrix entries processed. This comparison is only an
order-of-magnitude cost model: sparse matrix-vector products perform
a different mix of floating-point operations and memory accesses, and
ARPACK also incurs vector orthogonalization and restart overhead. The
measured speedup is therefore substantially smaller than the ratio of
these nominal operation counts.

\subsection{Numerical precision}
\label{sec:precision}

Table~\ref{tab:precision} compares the dense and sparse ground-state
energies and gives the Lanczos residual for the Heisenberg chain
($N \le 16$; larger sizes have no dense reference).

\begin{table}[t]
\centering
\small
\caption{Agreement between dense and sparse eigensolvers for the
Heisenberg chain. The relative error is
$|E_0^{\text{dense}} - E_0^{\text{sparse}}| / |E_0^{\text{dense}}|$.
The dense and sparse implementations agree with each other; they are
not compared against an exact mathematical reference.}
\label{tab:precision}
\begin{tabular}{rrrrr}
\toprule
$N$ & $D$ & $|E_0^{\text{dense}} - E_0^{\text{sparse}}|$
    & relative error & $\|H\psi_0 - E_0\psi_0\|$ \\
\midrule
4  & 6     & $1.8 \times 10^{-15}$ & $2.7 \times 10^{-16}$ & $2.4 \times 10^{-15}$ \\
6  & 20    & $5.3 \times 10^{-15}$ & $5.3 \times 10^{-16}$ & $8.0 \times 10^{-15}$ \\
8  & 70    & $1.1 \times 10^{-14}$ & $7.9 \times 10^{-16}$ & $1.5 \times 10^{-14}$ \\
10 & 252   & $1.8 \times 10^{-14}$ & $1.0 \times 10^{-15}$ & $1.0 \times 10^{-14}$ \\
12 & 924   & $3.2 \times 10^{-14}$ & $1.6 \times 10^{-15}$ & $1.4 \times 10^{-14}$ \\
14 & 3432  & $1.2 \times 10^{-13}$ & $5.2 \times 10^{-15}$ & $3.3 \times 10^{-14}$ \\
16 & 12870 & $9.6 \times 10^{-14}$ & $3.5 \times 10^{-15}$ & $2.2 \times 10^{-14}$ \\
\bottomrule
\end{tabular}
\end{table}

The two solvers agree with each other to a relative accuracy at or
below $5 \times 10^{-15}$, i.e.\ a few units of machine epsilon, and
the residual norms are $\lesssim 3.3 \times 10^{-14}$. For the tested
Heisenberg sizes, the choice of solver therefore has no measurable
effect on the reported ground-state energy. For a Hermitian matrix,
the residual provides a direct a posteriori bound on the distance
between an approximate eigenvalue and the spectrum: for a normalized
vector $\psi$ and $r = H\psi - \tilde E\psi$, there exists an exact
eigenvalue $\lambda$ such that
\begin{equation}
|\lambda - \tilde E| \le \|r\|.
\label{eq:hermitian-bound}
\end{equation}

\subsection{Symmetry-sector reduction}
\label{sec:symmetry}

\paragraph{Burnside counting for the Heisenberg chain.}
Beyond $S^z$, the open Heisenberg chain commutes with spatial
reflection $R$ and global spin flip $P$. For the even-$N$ chains
studied here, the ground state lies in the sector
$P = R = (-1)^{N/2}$, placing it in the $(+,+)$ sector for
$N \equiv 0 \pmod 4$ and in $(-,-)$ for $N \equiv 2 \pmod 4$. The
symmetry assignment is consistent with Marshall's sign
rule~\cite{Marshall1955,LiebMattis1962}.

Burnside's lemma for $\mathbb{Z}_2\times\mathbb{Z}_2$ gives
\begin{equation}
D_{(+,+)} = \tfrac14\left[
  D_{S^z} + \mathrm{tr}(R) + \mathrm{tr}(RP) \right],
\label{eq:burnside}
\end{equation}
because $\mathrm{tr}(P) = 0$ in the $S^z=0$ sector, with
$\mathrm{tr}(R) = \binom{N/2}{N/4}$ for $N/2$ even and $0$ for $N/2$
odd, and $\mathrm{tr}(RP) = 2^{N/2}$.

\begin{table}[t]
\centering
\small
\caption{Heisenberg-chain sector dimensions. Last two columns:
reduction relative to the full space and to the $S^z=0$ sector.}
\label{tab:symmetry}
\begin{tabular}{rrrrrr}
\toprule
$N$ & $D_{\text{full}}$ & $D_{S^z}$ & $D_{(+,+)}$
    & $D_{\text{full}}/D_{(+,+)}$ & $D_{S^z}/D_{(+,+)}$ \\
\midrule
4  & 16    & 6     & 3    & 5.33  & 2.00 \\
6  & 64    & 20    & 7    & 9.14  & 2.86 \\
8  & 256   & 70    & 23   & 11.13 & 3.04 \\
10 & 1024  & 252   & 71   & 14.42 & 3.55 \\
12 & 4096  & 924   & 252  & 16.25 & 3.67 \\
14 & 16384 & 3432  & 890  & 18.41 & 3.86 \\
16 & 65536 & 12870 & 3299 & 19.87 & 3.90 \\
\bottomrule
\end{tabular}
\end{table}

The factor $D_{S^z}/D_{(+,+)}$ rises from $2.0$ at $N=4$ to $3.90$ at
$N=16$ and tends to $4$ from below (Fig.~\ref{fig:hilbert}).

\begin{figure}[t]
\centering
\safeincludegraphics[width=0.85\textwidth]{Fig2_hilbert_scaling.pdf}
\caption{Hilbert-space dimension versus $N$ for the full space, the
$S^z=0$ sector, and the combined $S^z, R, P$ sector.}
\label{fig:hilbert}
\end{figure}

\paragraph{Parity sector of the TFI chain.}
$H_{\mathrm{TFI}}$ commutes with $P = \prod_i \sigma^x_i$. By
Perron--Frobenius, the ground state is unique with positive
amplitudes and hence $P=+1$. The $P=+1$ sector has dimension
$2^{N-1}$.

Table~\ref{tab:tfi-parity} shows that the reduced Hamiltonian
reproduces Eq.~\eqref{eq:tfi-closed} with maximum absolute deviation
$4.6 \times 10^{-14}$ for $N = 12$--$18$, i.e.\ to machine precision.

\begin{table}[t]
\centering
\small
\caption{TFI ground-state energy in the $P=+1$ sector versus the
closed form, Eq.~\eqref{eq:tfi-closed}.}
\label{tab:tfi-parity}
\begin{tabular}{rrrr}
\toprule
$N$ & $D_{P=+1}$ & $E_0$ (numerical) & $|E_0^{\mathrm{num}} - E_0^{\mathrm{cf}}|$ \\
\midrule
12 & 2048   & $-14.925971110$ & $4.6\times10^{-14}$ \\
14 & 8192   & $-17.471004055$ & $0$ \\
16 & 32768  & $-20.016387900$ & $2.1\times10^{-14}$ \\
17 & 65536  & $-21.289173773$ & $4.3\times10^{-14}$ \\
18 & 131072 & $-22.562008725$ & $2.5\times10^{-14}$ \\
\bottomrule
\end{tabular}
\end{table}

Table~\ref{tab:K-parity} reports the Lanczos iteration count in both
parity sectors; the iteration counts are deterministic given the
fixed starting vector and can be reproduced on any machine with the
same library versions. Table~\ref{tab:parity-timing} reports the
wall-clock comparison between the full space and the $P=+1$ sector
measured in the same logged session as the other timing tables.

\begin{table}[t]
\centering
\small
\caption{Lanczos iteration count for the TFI chain in the two parity
sectors (minimum over 10 repetitions). $K$ values are deterministic
given the fixed \texttt{v0}.}
\label{tab:K-parity}
\begin{tabular}{rrrr}
\toprule
$N$ & $D_{P=+1}$ & $K_{P=+1}$ & $K_{P=-1}$ \\
\midrule
10 & 512   & 71  & 71  \\
12 & 2048  & 81  & 91  \\
14 & 8192  & 91  & 101 \\
16 & 32768 & 111 & 121 \\
\bottomrule
\end{tabular}
\end{table}

\begin{table}[t]
\centering
\small
\caption{TFI chain, full space versus $P=+1$ sector. Measured in the
same logged session as the other timing tables. Minimum of 10; times
in seconds.}
\label{tab:parity-timing}
\begin{tabular}{rrrrrrr}
\toprule
$N$ & $D_{\mathrm{full}}$ & $D_{P=+1}$ & $t_{\mathrm{full}}$ & $t_{P=+1}$ & ratio & $K_{\mathrm{full}}/K_{P=+1}$ \\
\midrule
10 & 1024  & 512   & $0.0053$ & $0.0028$ & $1.89$ & 91/71 \\
12 & 4096  & 2048  & $0.0167$ & $0.0082$ & $2.05$ & 101/81 \\
14 & 16384 & 8192  & $0.0805$ & $0.0306$ & $2.63$ & 121/91 \\
16 & 65536 & 32768 & $0.7200$ & $0.1908$ & $3.77$ & 141/111 \\
\bottomrule
\end{tabular}
\end{table}

\begin{figure}[t]
\centering
\safeincludegraphics[width=0.85\textwidth]{Fig4_tfi_parity.pdf}
\caption{TFI chain in the full space and in the $P=+1$ sector:
Hilbert-space dimension and ground-state wall-clock time versus $N$.}
\label{fig:parity}
\end{figure}

The parity speedup is $1.9$--$3.8\times$. Halving $D$ alone would
give $2\times$; the product of the dimension ratio and the iteration
ratio $K_{\mathrm{full}}/K_{\mathrm{par}}$ predicts $1.3$, $1.2$,
$1.3$, $1.3$ for $N=10,12,14,16$. The measured ratios exceed these
predictions because the smaller matrix also fits more efficiently in
the cache, an effect not captured by the iteration-count model.

\subsection{Finite-size convergence}
\label{sec:convergence}

\paragraph{Open-chain expansion and the exact TFI case.}
For open chains
\begin{equation}
E_0(N) = N\epsilon_\infty + \sigma + \mathcal{O}(1/N),\qquad
\frac{E_0}{N} = \epsilon_\infty + \frac{\sigma}{N} + \mathcal{O}(1/N^2),
\label{eq:obc-expansion}
\end{equation}
with $\sigma$ the surface energy. For the critical TFI chain,
expanding Eq.~\eqref{eq:tfi-closed} gives
\begin{equation}
\frac{E_0}{N} = -\frac4\pi + \frac{1-2/\pi}{N} - \frac{\pi}{24N^2}
  + \frac{\pi}{48N^3}
  - \Bigl(\frac{\pi}{96} + \frac{7\pi^3}{23040}\Bigr)\frac{1}{N^4}
  + \mathcal{O}(N^{-5}),
\label{eq:tfi-expansion}
\end{equation}
so $\sigma_{\mathrm{TFI}} = 1-2/\pi \approx 0.36338$. The $N^{-2}$
coefficient equals $-\pi c v/24$ with $c=1/2$ and $v=2$ (Pauli
normalization), the open-boundary conformal-field-theory
form~\cite{Cardy1984}.

\begin{table}[t]
\centering
\small
\caption{TFI surface energy. $g(N) = (E_0/N-\epsilon_\infty) N$
approaches $\sigma_{\mathrm{TFI}} = 1-2/\pi$ from below; the residual
is the difference between $g(N)$ and $\sigma-\pi/(24N)$.}
\label{tab:surface}
\begin{tabular}{rrrrr}
\toprule
$N$ & $E_0/N$ & $g(N)$ & $\sigma - \pi/(24N)$ & residual \\
\midrule
4  & $-1.189693$ & $0.33419$ & $0.33065$ & $0.00354$ \\
6  & $-1.216038$ & $0.34321$ & $0.34156$ & $0.00165$ \\
8  & $-1.229744$ & $0.34796$ & $0.34702$ & $0.00094$ \\
10 & $-1.238149$ & $0.35091$ & $0.35029$ & $0.00062$ \\
12 & $-1.243831$ & $0.35290$ & $0.35247$ & $0.00043$ \\
14 & $-1.247929$ & $0.35435$ & $0.35404$ & $0.00031$ \\
16 & $-1.251024$ & $0.35544$ & $0.35520$ & $0.00024$ \\
$\infty$ & $-4/\pi$ & $0.36338$ & $0.36338$ & $0$ \\
\bottomrule
\end{tabular}
\end{table}

\paragraph{Extrapolation.}
Table~\ref{tab:extrap} compares four least-squares fits. At
$N_{\max} = 16$ the three-parameter OBC error is
$9.8\times10^{-5}$, the two-parameter PBC error is
$1.8\times10^{-2}$, and adding the $N^{-3}$ term reduces the OBC
error to $6.3\times10^{-6}$, a factor of $15$. The PBC form gives a
relative extrapolation error of $1.8\%$ because it lacks the surface
term altogether (Fig.~\ref{fig:extrap}).

\begin{table}[t]
\centering
\small
\caption{Relative error of the extrapolated bulk energy for the
critical TFI chain. The tabulated error is the extrapolation error to
$N=\infty$, not a fit residual. The four-parameter fit is not shown
at $N_{\max}=8$ because three data points do not constrain four
parameters.}
\label{tab:extrap}
\begin{tabular}{rrrrrr}
\toprule
$N_{\max}$ & points & 2-param OBC & 3-param OBC & 4-param OBC & 2-param PBC \\
\midrule
8  & 3 & $2.88\times10^{-3}$ & $1.94\times10^{-4}$ & ---                 & $2.51\times10^{-2}$ \\
10 & 4 & $2.44\times10^{-3}$ & $1.56\times10^{-4}$ & $1.12\times10^{-5}$ & $2.27\times10^{-2}$ \\
12 & 5 & $2.13\times10^{-3}$ & $1.31\times10^{-4}$ & $8.95\times10^{-6}$ & $2.09\times10^{-2}$ \\
14 & 6 & $1.89\times10^{-3}$ & $1.12\times10^{-4}$ & $7.40\times10^{-6}$ & $1.94\times10^{-2}$ \\
16 & 7 & $1.70\times10^{-3}$ & $9.77\times10^{-5}$ & $6.27\times10^{-6}$ & $1.81\times10^{-2}$ \\
\bottomrule
\end{tabular}
\end{table}

The three-parameter errors decrease only slowly with $N_{\max}$
because Eq.~\eqref{eq:tfi-expansion} contains a $+\pi/(48N^3)$ term
that the fit omits. A four-parameter fit including this term reduces
the extrapolation error at $N_{\max} = 16$ from $9.8\times10^{-5}$ to
$6.3\times10^{-6}$ (Table~\ref{tab:extrap}), a consistency check
rather than an independently achieved accuracy, because the added
term was predicted by the same analytic expansion under test. The
fitted $N^{-3}$ coefficient is $0.046$, about $30\%$ below the
analytic value $\pi/48 \approx 0.065$. Equation~\eqref{eq:tfi-expansion}
gives a negative $N^{-4}$ coefficient
$-\pi/96 - 7\pi^3/23040 \approx -0.042$; a truncated fit omitting
this term would tend to underestimate the $N^{-3}$ coefficient.
Pinning the $N^{-4}$ coefficient at its analytic value and refitting
the remaining terms raises the inferred $N^{-3}$ coefficient to
$\approx 0.071$. The free and pinned estimates therefore bracket
$\pi/48$ from below and above. The two together confirm that the
analytic expansion is consistent with the data, but the available
system sizes do not resolve the $N^{-3}$ coefficient to better than
about $30\%$. The three-parameter fit also biases the fitted
parameters (Table~\ref{tab:fitparams}): at $N_{\max} = 16$ it gives
$\sigma = 0.3602$ ($0.9\%$ below $0.36338$) and $\gamma = -0.1061$
($19\%$ smaller in magnitude than $-\pi/24 = -0.1309$).

\begin{table}[t]
\centering
\small
\caption{Parameters of the three- and four-parameter OBC fits. Exact
values: $\epsilon_\infty = -1.273240$, $\sigma = 0.363380$,
$\gamma = -0.130900$, $\pi/48 = 0.065450$. The four-parameter column
is not shown at $N_{\max}=8$ (underdetermined).}
\label{tab:fitparams}
\begin{tabular}{rrrrrr}
\toprule
$N_{\max}$ & points & $\epsilon_\infty$ (3p) & $\sigma$ (3p) & $\gamma$ (3p) & $\delta$ (4p) \\
\midrule
8  & 3 & $-1.272992$ & $0.358771$ & $-0.102295$ & --- \\
10 & 4 & $-1.273040$ & $0.359295$ & $-0.103620$ & $0.04480$ \\
12 & 5 & $-1.273073$ & $0.359676$ & $-0.104624$ & $0.04544$ \\
14 & 6 & $-1.273097$ & $0.359966$ & $-0.105416$ & $0.04594$ \\
16 & 7 & $-1.273115$ & $0.360196$ & $-0.106059$ & $0.04635$ \\
\bottomrule
\end{tabular}
\end{table}

\begin{figure}[t]
\centering
\safeincludegraphics[width=0.85\textwidth]{Fig3_extrapolation.pdf}
\caption{Extrapolation error versus the largest system size $N_{\max}$
used in the fit.}
\label{fig:extrap}
\end{figure}

\paragraph{Critical gap.}
The TFI gap satisfies $N\,\Delta E = 2.778,\ 2.893,\ 2.953,\ 2.989,\
3.014,\ 3.032,\ 3.045,\ 3.056$ for $N = 4,6,\dots,18$, approaching
$\pi$ from below (Fig.~\ref{fig:gap}), consistent with
$4\sin[\pi/(4N+2)]$.

\begin{figure}[t]
\centering
\safeincludegraphics[width=0.85\textwidth]{Fig6_gap_convergence.pdf}
\caption{Critical TFI gap multiplied by $N$ versus $N$; the limit is
$\pi$. The prediction is $4\sin[\pi/(4N+2)] \cdot N$, the free-fermion
single-particle gap.}
\label{fig:gap}
\end{figure}

\paragraph{Heisenberg chain.}
Table~\ref{tab:heis-surface} lists $E_0/N$ and
$g(N) = (E_0/N-\epsilon_\infty) N$ for the Heisenberg chain with
$\epsilon_\infty = 1 - 4\ln 2$. Two-point estimates of the
coefficient of the $1/N$ correction from $(4,8)$, $(8,12)$, $(12,16)$
give $\gamma \approx -0.44$, $-0.52$, and $-0.56$. The isotropic
Heisenberg chain contains logarithmic finite-size corrections from
the marginally irrelevant operator at the SU(2)-symmetric
point~\cite{Affleck1989}, so a fit involving a surface term,
subleading powers of $1/N$, and logarithmic corrections is not well
constrained by these sizes. We therefore do not quote a Heisenberg
surface energy.

\begin{table}[t]
\centering
\small
\caption{Heisenberg energy per site,
$\epsilon_\infty = 1 - 4\ln 2 \approx -1.772589$.}
\label{tab:heis-surface}
\begin{tabular}{rrr}
\toprule
$N$ & $E_0/N$ & $(E_0/N-\epsilon_\infty)N$ \\
\midrule
4  & $-1.616025$ & $0.626253$ \\
6  & $-1.662385$ & $0.661224$ \\
8  & $-1.687466$ & $0.680979$ \\
10 & $-1.703214$ & $0.693746$ \\
12 & $-1.714030$ & $0.702702$ \\
14 & $-1.721921$ & $0.709343$ \\
16 & $-1.727934$ & $0.714471$ \\
\bottomrule
\end{tabular}
\end{table}

\subsection{Entanglement entropy and central charge}
\label{sec:entropy}

The half-chain entropy of the Heisenberg ground state is computed
from the singular values of the ground-state vector reshaped into a
$2^{N/2} \times 2^{N/2}$ matrix,
$S = -\sum_i \sigma_i^2 \ln \sigma_i^2$, and is compared with the
open-chain form $S \simeq (c/6)\ln N + \mathrm{const}$
\cite{Calabrese2004}. Table~\ref{tab:entropy} shows a strong even-odd
structure in $\ell = N/2$ (Fig.~\ref{fig:ent}). Such alternating
contributions in open critical chains are
known~\cite{Laflorencie2006}.

\begin{table}[t]
\centering
\small
\caption{Half-chain entropy (natural logarithm) of the Heisenberg
chain for even $N$.}
\label{tab:entropy}
\begin{tabular}{rrr}
\toprule
$N$ & $N \bmod 4$ & $S$ \\
\midrule
4  & 0 & $0.319368$ \\
6  & 2 & $0.711373$ \\
8  & 0 & $0.456976$ \\
10 & 2 & $0.737869$ \\
12 & 0 & $0.536833$ \\
14 & 2 & $0.762250$ \\
16 & 0 & $0.592307$ \\
\bottomrule
\end{tabular}
\end{table}

A linear fit of $S$ against $\ln N$ over the four
$N \equiv 0 \pmod 4$ points gives slope $0.1971 \pm 0.0008$ and hence
$c_{\mathrm{eff}} = 1.183 \pm 0.005$ against $c=1$. The formal
uncertainty $\pm 0.005$ reflects only the scatter of the four points
around the fitted line. The $N \equiv 2 \pmod 4$ branch (three points)
gives a least-squares slope $0.059$, i.e.\ $c_{\mathrm{eff}} \approx
0.36$. The two branches therefore yield effective central charges
differing by a factor of three, and neither branch alone is a
reliable estimator at these sizes.

\begin{figure}[t]
\centering
\safeincludegraphics[width=0.85\textwidth]{Fig5_entanglement.pdf}
\caption{Half-chain entanglement entropy of the Heisenberg chain
versus $\ln N$. The parity effect splits the data into two branches.}
\label{fig:ent}
\end{figure}

% =========================================================================
\section{Discussion}
\label{sec:discussion}

\subsection{Recommendations}
\label{sec:recs}

For ground-state or few-state calculations, our measurements suggest
that sparse Lanczos becomes preferable at $D$ somewhat above the
measured break-even scale. The measured crossover is
$D \approx 2\times10^{2}$ with the numpy baseline. Because
\texttt{scipy.linalg.eigvalsh} is faster by a factor of $2.0$--$2.5$
at large $D$, its crossover with the sparse solver would lie higher;
this was not measured directly. We recommend a working threshold of
$D \gtrsim 3\times10^{2}$, which is safe under either library choice.
Dense diagonalization remains the natural choice when the full
spectrum or all eigenvectors are needed, or for $D \lesssim 10^{2}$
where the cost is negligible. The benefit becomes substantial where
it matters: about $3.4\times10^{2}$ at $D \approx 2000$ and
$1.1\times10^{3}$ at $D = 4096$ on our hardware, using the numpy
dense baseline.

A practical reduction in runtime is obtained by relaxing the Lanczos
tolerance from the tightest ARPACK setting to \texttt{tol} $=
10^{-10}$; in the tested TFI cases this reduced the runtime by
$20$--$40\%$ while maintaining converged ground-state energies.

For open chains the fit ansatz must contain the $1/N$ surface term;
the TFI chain provides a free exact benchmark. Symmetry reduction by
$S^z$ should always be applied; reflection and spin flip add a further
factor approaching $4$ for a known ground-state sector, and the TFI
parity sector gives a $1.9$--$3.8\times$ speedup. Central charges
should not be extracted from a single $N \bmod 4$ branch of $S_{N/2}$
at $N \le 16$.

\subsection{Limitations}
\label{sec:limitations}

All timings come from one machine, one single-threaded BLAS build, and
one SciPy/ARPACK version; the ratios quoted are specific to them. The
calibration session was conducted on a single day; variation across
days, thermal loads, and background processes was not measured and
may be larger. The dense--sparse comparison is not strictly
apples-to-apples, since the dense solver returns the entire spectrum
while the sparse solver returns only the ground state; restricting
the dense computation to the lowest eigenvalue changes the timing by
$\approx 2$--$13\%$ at $D \geq 924$ (Table~\ref{tab:fair}), so the
reported ratios are not appreciably biased by the spectral scope of
the dense call. We did not benchmark against QuSpin or other
established packages; the scope of this study is the behaviour of
the default NumPy/SciPy ED workflow, and adding a packaged library
would change the comparison to a different question. The reported
hardware lacks AVX512, which affects the dense timings in
particular; the effective numpy throughput at $D=4096$ is
$3.3$~GFLOP/s, well below the peak of the processor. On a
multi-threaded machine with AVX512 the dense baseline would be faster
and the dense-to-sparse ratios correspondingly smaller. The purpose
of the benchmark is not to establish universal hardware-independent
crossover constants, but to quantify the scale of effects that can
arise in a reproducible small-scale ED workflow.

% =========================================================================
\section{Conclusions}
\label{sec:conclusions}

\begin{enumerate}
  \item For the single-threaded NumPy/SciPy implementation tested
        here, the dense-to-sparse ground-state workflow crossover
        occurs at $D \approx 2\times10^{2}$ with the
        \texttt{numpy.linalg.eigvalsh} dense baseline. At
        $D = 4096$ the workflow ratio is $\approx 1.1\times10^{3}$
        with that baseline and $\approx 6.6\times10^{2}$ with
        \texttt{scipy.linalg.eigvalsh}; the two libraries differ by
        a factor of $2.3$ on the same matrix. Restricting the dense
        call to the lowest eigenvalue provides an additional
        $\approx 2$--$13\%$ saving at $D \geq 924$.

  \item The measured Lanczos iteration counts, $K = 71$--$141$ for
        $N = 8$--$16$, are substantially larger than a commonly
        assumed value $K \approx 30$.

  \item For the tested Heisenberg sizes, dense and sparse
        ground-state energies agree to relative differences below
        $5.2\times10^{-15}$, with sparse residual norms below
        $3.3\times10^{-14}$.

  \item For open chains, the leading finite-size correction to
        $E_0/N$ is a surface term proportional to $1/N$. An
        OBC-consistent three-parameter fit recovers the TFI bulk
        energy to approximately $10^{-4}$ at $N_{\max} = 16$;
        adding the $N^{-3}$ term predicted by the analytic
        expansion reduces the fit residual to $6.3\times10^{-6}$.
        This residual is a consistency check rather than an
        independently achieved accuracy, because the added term is
        predicted by the same analytic expansion under test. The
        fitted coefficient is $30\%$ below $\pi/48$; pinning the
        analytic $N^{-4}$ coefficient in the fit moves it to $9\%$
        above, so the two estimates bracket the analytic value and
        confirm the consistency of the expansion, but do not
        resolve the $N^{-3}$ coefficient to better than about
        $30\%$ at the available sizes.

  \item Symmetry reduction provides a substantial reduction in the
        Hilbert-space dimension. In the Heisenberg chain, restricting
        to the appropriate $S^z$, reflection, and spin-flip sector
        reduces the dimension by a factor growing from $5.3$ at
        $N=4$ to $19.9$ at $N=16$. For the TFI chain, restriction to
        the ground-state parity sector gives a factor-of-two
        reduction in dimension and measured speedups of
        $1.9$--$3.8\times$ in the tested range.

  \item Half-chain entanglement entropies of the open Heisenberg
        chain show a strong $N \bmod 4$ dependence at $N \le 16$.
        Single-branch fits on three or four points give effective
        central charges differing by a factor of three between the
        two branches, demonstrating that neither branch alone is a
        reliable estimator at these sizes.
\end{enumerate}

% =========================================================================
\section*{Declarations}

\subsection*{Funding}
The authors received no specific funding for this work.

\subsection*{Competing interests}
The authors declare no competing interests.

\subsection*{Ethics approval}
Not applicable.

\subsection*{Author contributions}
S.S.\ designed the study, wrote the code, performed the numerical
experiments, analyzed the data, and wrote the manuscript. M.I.J.\
contributed to the analysis of the numerical results and to the
revision of the manuscript. Both authors read and approved the final
manuscript.

\subsection*{Data availability}
The Python module, benchmark scripts, raw CSV data, figure and table
generators, and the log from the calibration session are available
in the
\href{https://github.com/sswaid991-bit/ed-spin-chains}{public
repository}. The repository includes the logged session
\texttt{20261007\_070101} with the exact library files and OpenBLAS
versions in use, and the per-session CSVs for every table in this
manuscript. Rerunning \texttt{reproduce.py} reproduces the full set
of tables and figures in about 25 minutes on the reference laptop.

% =========================================================================
\begin{thebibliography}{99}

\bibitem{Hulthen1938}
L.~Hulth\'en, \"Uber das Austauschproblem eines Kristalles, Ark.\
Mat.\ Astron.\ Fys.\ \textbf{26A}, 1 (1938)

\bibitem{Pfeuty1970}
P.~Pfeuty, The one-dimensional Ising model with a transverse field,
Ann.\ Phys.\ \textbf{57}, 79 (1970).
\href{https://doi.org/10.1016/0003-4916(70)90270-8}{https://doi.org/10.1016/0003-4916(70)90270-8}

\bibitem{Sandvik2010}
A.~W.~Sandvik, Computational studies of quantum spin systems, AIP
Conf.\ Proc.\ \textbf{1297}, 135 (2010).
\href{https://doi.org/10.1063/1.3518900}{https://doi.org/10.1063/1.3518900}

\bibitem{Weinberg2017}
P.~Weinberg, M.~Bukov, QuSpin: a Python package for dynamics and exact
diagonalisation of quantum many-body systems, SciPost Phys.\
\textbf{2}, 003 (2017).
\href{https://doi.org/10.21468/SciPostPhys.2.1.003}{https://doi.org/10.21468/SciPostPhys.2.1.003}

\bibitem{Schollwock2011}
U.~Schollw\"ock, The density-matrix renormalization group in the age
of matrix product states, Ann.\ Phys.\ \textbf{326}, 96 (2011).
\href{https://doi.org/10.1016/j.aop.2010.09.012}{https://doi.org/10.1016/j.aop.2010.09.012}

\bibitem{Lehoucq1998}
R.~B.~Lehoucq, D.~C.~Sorensen, C.~Yang, \emph{ARPACK Users' Guide}
(SIAM, Philadelphia, 1998).
\href{https://doi.org/10.1137/1.9780898719628}{https://doi.org/10.1137/1.9780898719628}

\bibitem{Marshall1955}
W.~Marshall, Antiferromagnetism, Proc.\ R.\ Soc.\ London A
\textbf{232}, 48 (1955).
\href{https://doi.org/10.1098/rspa.1955.0204}{https://doi.org/10.1098/rspa.1955.0204}

\bibitem{LiebMattis1962}
E.~H.~Lieb, D.~Mattis, Ordering energy and specific heat of a linear
antiferromagnet, J.\ Math.\ Phys.\ \textbf{3}, 749 (1962).
\href{https://doi.org/10.1063/1.1724280}{https://doi.org/10.1063/1.1724280}

\bibitem{Cardy1984}
J.~L.~Cardy, Conformal invariance and surface critical behavior,
Nucl.\ Phys.\ B \textbf{240}, 514 (1984).
\href{https://doi.org/10.1016/0550-3213(84)90241-4}{https://doi.org/10.1016/0550-3213(84)90241-4}

\bibitem{Affleck1989}
I.~Affleck, D.~Gepner, H.~J.~Schulz, T.~Ziman, Critical behaviour of
spin-$s$ Heisenberg antiferromagnetic chains: analytic and numerical
results, J.\ Phys.\ A \textbf{22}, 511 (1989).
\href{https://doi.org/10.1088/0305-4470/22/5/005}{https://doi.org/10.1088/0305-4470/22/5/005}

\bibitem{Calabrese2004}
P.~Calabrese, J.~Cardy, Entanglement entropy and quantum field
theory, J.\ Stat.\ Mech.\ (2004) P06002.
\href{https://doi.org/10.1088/1742-5468/2004/06/P06002}{https://doi.org/10.1088/1742-5468/2004/06/P06002}

\bibitem{Laflorencie2006}
N.~Laflorencie, E.~S.~S{\o}rensen, M.-S.~Chang, I.~Affleck, Boundary
effects in the critical scaling of entanglement entropy in 1D
systems, Phys.\ Rev.\ Lett.\ \textbf{96}, 100603 (2006).
\href{https://doi.org/10.1103/PhysRevLett.96.100603}{https://doi.org/10.1103/PhysRevLett.96.100603}

\end{thebibliography}

% =========================================================================
\appendix

\section{XXZ reference data}
\label{app:xxz}

\begin{table}[t]
\centering
\small
\caption{Ground-state energy per site of the four models (Pauli
convention).}
\label{tab:xxz}
\begin{tabular}{rrrrr}
\toprule
$N$ & TFI & Heisenberg & XXZ ($\Delta=0.5$) & XXZ ($\Delta=2.0$) \\
\midrule
4  & $-1.189693$ & $-1.616025$ & $-1.356086$ & $-2.186141$ \\
6  & $-1.216038$ & $-1.662385$ & $-1.398592$ & $-2.262953$ \\
8  & $-1.229744$ & $-1.687466$ & $-1.421625$ & $-2.305637$ \\
10 & $-1.238149$ & $-1.703214$ & $-1.436100$ & $-2.333073$ \\
12 & $-1.243831$ & $-1.714030$ & $-1.446049$ & $-2.352318$ \\
14 & $-1.247929$ & $-1.721921$ & $-1.453312$ & $-2.366633$ \\
16 & $-1.251024$ & $-1.727934$ & $-1.458847$ & $-2.377738$ \\
\bottomrule
\end{tabular}
\end{table}

\end{document}